\documentclass[11pt]{article}
\usepackage[margin=1in]{geometry}
\usepackage[T1]{fontenc}
\usepackage{newtxtext}
\usepackage{microtype}
\usepackage{amsmath,amsthm,mathtools}
\usepackage{newtxmath}
\usepackage{booktabs,array,longtable}
\usepackage{enumitem}
\usepackage{xcolor}
\usepackage{listings}
\usepackage{hyperref}
\usepackage{fancyhdr}
\hypersetup{colorlinks=true,linkcolor=blue!45!black,citecolor=blue!45!black,urlcolor=blue!45!black,
 pdftitle={Sparse Certified Port Incidence for Compressed Unknot-Recognition Pipelines},
 pdfauthor={ChatGPT; research continuation for ProveIt}}
\pagestyle{fancy}
\fancyhf{}
\fancyhead[L]{\small Sparse certified port incidence}
\fancyhead[R]{\small ProveIt research continuation}
\fancyfoot[C]{\thepage}
\setlength{\headheight}{14pt}
\setlength{\emergencystretch}{2em}
\setlist{itemsep=3pt,topsep=5pt}
\lstset{basicstyle=\ttfamily\small,breaklines=true,columns=fullflexible,keepspaces=true,
 frame=single,rulecolor=\color{black!25},aboveskip=9pt,belowskip=9pt}
\newtheorem{theorem}{Theorem}[section]
\newtheorem{lemma}[theorem]{Lemma}
\newtheorem{proposition}[theorem]{Proposition}
\newtheorem{corollary}[theorem]{Corollary}
\theoremstyle{definition}
\newtheorem{definition}[theorem]{Definition}
\newtheorem{example}[theorem]{Example}
\newtheorem{remark}[theorem]{Remark}
\newcommand{\supp}{\operatorname{supp}}
\newcommand{\poly}{\operatorname{poly}}
\newcommand{\bits}{\operatorname{bits}}
\newcommand{\univ}{[r]}
\newcommand{\NN}{\mathbb{N}}
\newcommand{\ZZ}{\mathbb{Z}}
\newcommand{\hist}{h}
\newcommand{\ind}{\mathbf{1}}
\newcommand{\emptysetmark}{\varnothing}
\title{\textbf{Sparse Certified Port Incidence for\protect\\
Compressed Unknot-Recognition Pipelines}\\[6pt]
\large Output-sensitive reconstruction, polynomial support bounds,\protect\\
and exact coning continuations}
\author{Research continuation for the ProveIt project\\
\small Prepared by ChatGPT for Vladimir Reshetnikov}
\date{8 October 2026}
\begin{document}
\maketitle
\begin{abstract}
The maintained marked-orbit interface in ProveIt reconstructs component
incidence with $r$ ports by a dense Boolean inversion with $2^r$ output slots.
We replace this representation with exact sparse recovery from the same
unweighted orbit-count primitive. For $s$ occurring signatures, balanced
block deletion uses
$O(s+\sum_{\varnothing\neq T\in\supp h}|T|\log(1+r/|T|))$ logical queries, without
knowing $s$ in advance. An independent certificate needs at most
$s+\sum_T|T|$ authenticated subset counts, plus the baseline count.

For interval-encoded ports, sparsity need not be assumed: a run-accounting
consequence of the classical weighted Agol--Hass--Thurston algorithm bounds
$s$ by $O(M+(k+1)^3(B+1))$, where $k$ is the number of pairings, $M$ the
number of marked intervals, and $B$ the endpoint bit length. Thus the new
black-box reduction is polynomial in its local encoded input, with no
logarithmic restriction on the number of ports. Polynomial weighted orbit
computation itself is classical, not a novelty claim of this note.

We also prove exact updates under coning and union re-marking, and refine
incidence by a supplied binary orientation character using only one further
query per occurring signature. Executable standard-library Python, adversarial
checks, an opt-in maintained-kernel adapter, and controlled component
measurements accompany the proofs. Sparse examples improve substantially;
a full-support control slows down. The measurements use a clearly labeled
audit extraction, not the complete maintained recognizer. These results
remove a local exponential representation bottleneck; they do not establish
an end-to-end quasi-polynomial unknot-recognition algorithm.
\end{abstract}
\vspace{4pt}
\noindent\textbf{Status.} Mathematical arguments and executable research artifacts;
not a formally verified or peer-reviewed theorem package. No repository writes
or production-default changes are included.
\newpage
\setcounter{tocdepth}{1}
\tableofcontents
\newpage

\section{The bottleneck and the scope of the improvement}
\label{sec:scope}

\subsection{What is being continued}
This note continues the interval-orbit branch of
\texttt{Topology/UnknotRecognition}, rather than restarting the project's
Khovanov, braid, determinant, or group-presentation work. The inspected
integration commit for the dense incidence module is
\texttt{0441aad3967b69621edabfdb85367d465b4fb32f}.
This is a baseline integration commit, not an assertion about the repository's
latest overall head. The exact inspected file hashes are recorded in the
package's \texttt{provenance.json}.

The relevant maintained article, \texttt{synthesis/orbit\_certificates.tex},
describes independently replayable unweighted orbit counts, coning a union of
marked intervals, exact caching of equal physical unions, and a dense
incidence histogram obtained by Boolean M\"obius inversion
\cite{repo-incidence,repo-article}. Query reuse reduces repeated work but
leaves an array of length $2^r$. Its conditional polynomial and
quasi-polynomial bounds consequently restrict the number of ports. The
interface reports properties of supplied interval relations, not a knot
verdict; ordinary recognition does not currently dispatch through it
\cite{repo-article}.

The precise target here is narrower and more actionable than replacing the
whole recognizer: retain the existing count/proof boundary, but eliminate
unnecessary enumeration of absent incidence signatures. This separation
matters. A local theorem can be fully unconditional in its encoded input
while its relevance to a global algorithm still depends on how those inputs
are generated.

\subsection{Established results, new deductions, and integration work}
The classical ingredients are the unweighted and weighted orbit algorithms
of Agol, Hass, and Thurston (AHT) \cite{aht}. Their Theorem 12 supplies a
polynomial unweighted count. Their weighted transfer lemma and Theorem 16
supply polynomial compressed orbit-weight data. We do not claim the first
polynomial algorithm for marked or weighted components.

The contributions developed here are the following precise combination.
First, a nonnegative sparse inversion can reuse only unweighted coning
queries, including their existing independent proofs. Second, short
minimal-support certificates avoid both dense inversion and replay of the
search schedule. Third, an explicit accounting argument connects interval
encoding to polynomially many occurring signatures. Fourth, a restricted
continuation language and a known-support double-cover calculation give
additional exact operations on the recovered representation. These are
proved below; priority beyond the inspected repository is not asserted.

The supplied implementation has two distinct assurance levels. Its generic
algebra, interval oracle, audit kernel, continuation operations, and signed
refinement have executed tests. The adapter for the maintained modules has
executed contract tests with independent literal witnesses. A genuine
maintained AHT trace run is supplied as a separate integration script but
was not executed in this delivery environment. No upstream full-suite pass
is attributed to this package.

\subsection{Relation to the quasi-polynomial goal}
Lackenby's July 2026 preprint develops incompressible-surface and hierarchy
methods relevant to unknot recognition. Section 9 explicitly separates the
number of hierarchy operations from the still-unspecified costs of some
operations, and leaves hierarchy-length and pattern-complexity parameters
unbounded there \cite{lackenby2026}. The present note addresses neither of
those global bounds. It identifies a class of aggregate component operations
whose representation need not contribute an additional $2^r$ factor.

An improvement from $2^r\poly(\text{local input})$ to
$\poly(\text{local input})$ is useful even before a hierarchy theorem is
complete. It changes which intermediate interfaces are computationally
plausible. It does not show that every relevant interface is an incidence
histogram, that all future attachments preserve this information, or that
a complete hierarchy has sufficiently small encoded descriptions.

\section{Interval relations and the coning oracle}
\label{sec:model}

\subsection{Input conventions}
Let $X_N=\{0,1,\ldots,N-1\}$. A pairing is an isometry between two nonempty
inclusive integer intervals of the same length. It is either a translation
or an order reversal. The smallest equivalence relation containing all
pairing arrows is denoted by $\sim$. Its classes are called \emph{orbits}.
There are $k$ supplied pairings and $C=|X_N/{\sim}|$ orbits.

Ports $A_1,\ldots,A_r$ are finite unions of half-open intervals
$[a,b)\cap\ZZ$, with $0\leq a\leq b\leq N$. Let $M$ be the number of
supplied marked interval records, before optional canonical merging. Empty,
overlapping, coincident, and nested ports are permitted. The convention
change is intentional: pairings use inclusive endpoints, while port records
use half-open endpoints. All endpoints and multiplicities are binary integers.
Put
\[
 B=\lceil\log_2(N+1)\rceil.
\]
The empty universe is dealt with directly. In polynomial expressions,
$B+1$ avoids degenerate zero factors.

For an orbit $O$, define its signature and the associated histogram by
\begin{align}
 T(O)&=\{i\in[r]:O\cap A_i\neq\varnothing\},\label{eq:signature}\\
 h(T)&=|\{O:T(O)=T\}|.\label{eq:histogram}
\end{align}
Write $s=|\supp h|$. The empty signature is retained. In particular, a
large region missed by all ports is not silently discarded. The requested
output is a list of the positive pairs $(T,h(T))$, not an array with a slot
for every subset of $[r]$.

An $r$-bit mask occupies $O(r)$ bits, even if its numerical value is close
to $2^r$. Merely using masks does not allocate a truth table. Nevertheless,
operations on those masks must be charged in the bit model rather than
regarded as constant time for unbounded $r$.

\subsection{Coning a physical union}
For a nonempty union $A$ of disjoint half-open intervals, add the following
relations. In each interval $[a,b)$ of length at least two, add the
translation $[a,b-2]\to[a+1,b-1]$. Choose one interval's first point as an
anchor and join every other interval's first point to it by a singleton
pairing. This construction uses at most twice the number of intervals in
$A$ and never expands an interval into points.

\begin{lemma}[Coning identity]
\label{lem:cone}
Let $C_A$ be the number of orbits after these relations are added, and let
$H(A)$ be the number of original orbits meeting $A$. Then
\[
 H(A)=C-C_A+1.
\]
For $A=\varnothing$, $H(A)=0$ without an auxiliary query.
\end{lemma}
\begin{proof}
Every point of $A$ is identified with the anchor. Hence every original orbit
meeting $A$ belongs to one new orbit. No orbit disjoint from $A$ is incident
to an added arrow, so it remains unchanged. Exactly $H(A)$ original orbits
are replaced by one, which gives the formula. This argument applies to any
initial equivalence relation; interval geometry is used only to encode the
extra arrows compactly.
\end{proof}

The identity and its interval construction are already present in the
maintained module \cite{repo-incidence}. We use them as an authenticated
oracle rather than changing the orbit scheduler.

For $U\subseteq[r]$, define
\begin{equation}
 G(U)=C-H\left(\bigcup_{i\notin U}A_i\right).
 \label{eq:zeta}
\end{equation}
An orbit contributes precisely when its signature is contained in $U$.
Thus
\begin{equation}
 G(U)=\sum_{T\subseteq U}h(T),\qquad G([r])=C.
 \label{eq:zeta-sum}
\end{equation}
When the physical complement-port union is nonempty,
$G(U)=C_A-1$; otherwise $G(U)=C$. A nonempty set of port indices may still
specify an empty physical union. The implementation tests the latter,
not merely whether the index mask is nonzero.

\subsection{What may be cached}
Canonicalize a physical union by sorting intervals and merging both overlap
and adjacency. Equal canonical unions lead to the same extra pairings and
the same source-bound count query. They may share a count and a proof.
Equality of numerical answers alone is not a valid cache key: two different
unions can have the same touched-orbit count but behave differently in other
queries. In particular, the sparse reduction and physical-union caching
are separate optimizations. The former removes absent output types; the
latter removes duplicate oracle inputs.

\section{Exact sparse inversion by positive residuals}
\label{sec:peeling}
This section applies to any nonnegative integer function on $2^{[r]}$,
not just to interval systems. Its hypothesis is the full oracle contract
\eqref{eq:zeta-sum}. Checking a few queried values lie between zero and $C$
does not establish that contract for an arbitrary untrusted oracle.

Suppose the algorithm has already recovered exact pairs
$(T_1,w_1),\ldots,(T_j,w_j)$. Define
\begin{equation}
 R_j(U)=G(U)-\sum_{\substack{1\leq a\leq j\\T_a\subseteq U}}w_a.
 \label{eq:residual}
\end{equation}
If the previous extractions are correct, this is the zeta transform of the
remaining nonnegative coefficients. Consequently it is monotone under
inclusion and nonnegative on every subset.

\begin{lemma}[Minimal positive support]
\label{lem:minimal}
If $R_j(U)>0$ and $R_j(U\setminus\{i\})=0$ for each $i\in U$, then $U$ is
an unrecovered signature and $R_j(U)=h(U)$.
\end{lemma}
\begin{proof}
Any proper subset $S\subsetneq U$ omits some $i\in U$. Its residual
coefficient is a nonnegative summand of $R_j(U\setminus\{i\})$, so it is
zero. The only possible positive residual summand in $R_j(U)$ is therefore
the coefficient at $U$ itself. Positivity proves that coefficient is present;
the displayed value is its entire multiplicity.
\end{proof}

A simple extraction begins with $U=[r]$ and tries deleting each index in
turn. A deletion is accepted exactly when the resulting residual remains
positive. A failed deletion remains impossible after later deletions:
if the candidate residual was zero, every smaller candidate also has zero
residual. At the end, Lemma~\ref{lem:minimal} applies. Importantly, one
extraction removes the full multiplicity of a type, not one component.
A coefficient equal to $2^{16000}$ therefore takes one extraction.

\begin{theorem}[Sparse recovery]
\label{thm:recovery}
Given $C$ and an exact oracle satisfying \eqref{eq:zeta-sum}, all positive
coefficients of $h$ can be recovered without knowing $s$ in advance.
Individual deletion uses at most $sr$ oracle requests. The procedure stores
only recovered entries and its query cache; it needs no $2^r$ array.
\end{theorem}
\begin{proof}
When the remaining mass $C-\sum_{a\leq j}w_a$ is positive, the starting
set $[r]$ has positive residual. Deletions preserve positivity and terminate
at a set satisfying Lemma~\ref{lem:minimal}. This recovers one previously
absent positive coefficient exactly and preserves the nonnegative-residual
invariant. There can be at most $s$ such extractions. Conversely, if any
coefficient remained positive, the remaining mass would be positive.
Termination at mass zero therefore proves completeness. There are at most
$r$ deletion probes per extraction. For $r=0$, the only possible coefficient
is $h(\varnothing)=C$ and no query is needed.
\end{proof}

A resource cap does not alter the theorem's acceptance condition. Exhausting
a cap before mass zero means that the complete output has not been obtained.
The adapter consequently returns \texttt{INCONCLUSIVE}, not a partial
histogram labeled complete and not a negative knot certificate.

\subsection{Why nonnegativity is essential}
For signed coefficients, cancellation can make $R(U)=0$ while positive and
negative proper-subset coefficients remain. A deletion can then remove
hidden support. For example, coefficients $h(\varnothing)=1$ and
$h(\{1\})=-1$ have total zero but nonempty support. The mass-zero stopping
rule would be invalid. Orbit multiplicities are nonnegative by definition,
which supplies exactly the structure used here.

\section{Balanced block deletion and its query complexity}
\label{sec:blocks}

\subsection{The algorithm}
Individual deletion spends a probe on every port even when a surviving type
uses very few ports. Instead, arrange $[r]$ in a balanced binary partition
tree. At a node with block $D$, query $R(U\setminus D)$. If it is positive,
delete the entire block and do not visit its descendants. If it is zero,
retain the block provisionally and visit both children. A singleton with a
zero deletion residual is kept. The traversal is deterministic and iterative;
its depth does not become a Python recursion requirement.

The algorithm's terminal set is still inclusion-minimal positive. A retained
singleton failed deletion when tested, and subsequent deletions cannot turn
a zero residual into a positive one. Deleted blocks have preserved the
existence of a surviving coefficient throughout.

\begin{theorem}[Output-sensitive block bound]
\label{thm:block}
Let one extraction return $T$, with $d=|T|$, and put
$b_r=\lceil\log_2 r\rceil$ for $r\geq1$. Its number of deletion probes is at most
\begin{equation}
 q_r(d)=1+2\sum_{\ell=0}^{b_r-1}\min(2^\ell,d).
 \label{eq:qexact}
\end{equation}
In particular, $q_r(0)=1$, and for $d>0$ it is
$O(1+d\log(1+r/d))$. The entire reconstruction uses at most
\begin{equation}
 Q\leq\sum_{T\in\supp h}q_r(|T|)
 \label{eq:qsum}
\end{equation}
logical requests. For $r=0$, set $Q=0$.
\end{theorem}
\begin{proof}
Consider a failed block $D$ tested at a current set $U_D$.
If the final support $T$ missed $D$, then
$T\subseteq U_D\setminus D$. The coefficient at $T$ belongs to the current
extraction's fixed residual and is positive, contradicting the failed test.
Every failed block thus intersects $T$.

At depth $\ell$, blocks are pairwise disjoint, so at most $d$ of them can
fail. There are also at most $2^\ell$ blocks at that depth. Only children
of failed internal nodes are queried. Counting the root probe and at most
two child probes per failed internal node gives \eqref{eq:qexact}, because
the balanced tree has height at most $b_r$. For $d=0$ the root succeeds and
there are no descendants. For $d>0$, split the sum where $2^\ell$ first
exceeds $d$. The early geometric sum is $O(d)$ and the remaining
$O(\log(1+r/d))$ levels contribute at most $d$ each. Summing over
extractions proves \eqref{eq:qsum}.
\end{proof}

Physical-union caching may make the number of actual orbit counts much
smaller than $Q$. Conversely, producing independent proofs incurs further
queries in the supplied two-pass implementation. These costs must not be
conflated with the logical discovery bound.

\subsection{A matching order of growth for a basic special case}
This search is closely related to group testing. The literature on binary
splitting and hypergraph reconstruction contains much broader query-model
results \cite{price-scarlett,bshouty-mazzawi}; we do not claim binary splitting
as a new general technique. The following elementary comparison explains
why the dependence on the size of an individual signature is reasonable.

\begin{proposition}[Singleton-profile lower bounds]
\label{prop:lower}
Consider deterministic algorithms correct for all singleton profiles
$h(T)=1$, with the unknown $T\subseteq[r]$. Among supports of size $d$,
some input needs at least $\lceil\log_2\binom rd\rceil$ oracle calls.
Moreover, for every nonempty actual $T$, its transcript needs at least
$|T|$ calls if correctness is required against supports of every size.
Thus the order $d\log(1+r/d)$ is optimal in the worst case at each $d>0$,
up to absolute constants, for this special case without a supplied rank.
\end{proposition}
\begin{proof}
Here $G(U)$ is one exactly when $T\subseteq U$; every answer is a bit.
Distinguishing $\binom rd$ possibilities gives the first decision-tree bound.
For the second, compare the actual $T$ with $T\setminus\{i\}$ for every
$i\in T$. A query distinguishes this pair only when $U$ contains all of
$T\setminus\{i\}$ and omits $i$. One query cannot have that property for
two different members of $T$. Every deletion alternative must be
excluded, requiring at least $|T|$ calls. If $d\leq r/2$, the binomial
bound is $\Omega(d\log(r/d))$; if $d>r/2$, the linear bound supplies
$\Omega(d\log(1+r/d))$.
\end{proof}

The qualification about a supplied rank is necessary. Knowing in advance
that $d=r$ determines $T$ without any query. We also make no claim that
summing the singleton lower bounds proves optimality for arbitrary
multi-support profiles; that is a different adaptive reconstruction problem.

\section{Short independent certificates}
\label{sec:certificates}

\subsection{Certificate data and algebraic checks}
A certificate lists positive pairs $(T_1,w_1),\ldots,(T_s,w_s)$ in an
extraction order. For each listed $T_j$, authenticate the following values:
\begin{equation}
 G(T_j)\quad\text{and}\quad G(T_j\setminus\{i\})\quad(i\in T_j).
 \label{eq:certificate-queries}
\end{equation}
Each authentication is either the baseline identity for an empty physical
union, or an independently checked coned orbit count. Repeated physical
unions share one proof. The verifier reconstructs all unions from the
original ports and binds each proof to the exact original pairing list plus
its deterministic coning rows.

At entry $j$, subtract the weights of previously listed types contained in
each query mask. Require residual $w_j$ at $T_j$ and zero at every tested
one-element deletion. Check masks are in range, types are distinct,
weights are positive integers rather than Boolean values, and
$\sum_jw_j=C$. The baseline count itself must be authenticated; it is not
merely a producer's declaration.

\begin{theorem}[Soundness and completeness of sparse replay]
\label{thm:cert}
Assume the underlying count proofs are sound for their bound inputs.
The checks above accept exactly complete, correct incidence profiles
presented in a valid minimal-extraction order. Every profile produced by
Theorem~\ref{thm:recovery} has such a certificate. The number of required
subset values before reuse is at most
\begin{equation}
 V\leq s+\sum_{T\in\supp h}|T|.
 \label{eq:vbound}
\end{equation}
The verifier need not call the discovery algorithm, reproduce its block
schedule, or perform a dense inverse transform.
\end{theorem}
\begin{proof}
Induct on the entry number. Before the first entry, all true coefficients
are nonnegative. If all preceding entries are exact coefficients, subtracting
them leaves the zeta transform of a nonnegative residual profile.
The authenticated checks for entry $j$ satisfy
Lemma~\ref{lem:minimal}, so the next entry is also exact. After all entries,
the residual has total mass zero; nonnegativity forces every omitted
coefficient to vanish. This proves soundness. Discovery produces precisely
these minimal positive residuals, so its entries satisfy the tests. Counting
one query for each type and one for each of its members gives
\eqref{eq:vbound}. None of these arguments relies on the producer's
particular deletion schedule.
\end{proof}

\subsection{Source binding and transport}
The sparse certificate is not a bare list of numerical oracle answers.
The maintained adapter stores one baseline local orbit proof, the sparse
payload, a table of local coned proofs, and a sparse list mapping verification
masks to proof indices. The checker rejects missing masks, repeated masks,
unused proofs, a proof index assigned to two different physical unions,
and distinct proof indices redundantly assigned to one physical union.
The native local checker is imported from
\texttt{fastunknot.interval\_orbit\_verify}, not from the producer.

The package includes exact tagged-hexadecimal JSON transport. It encodes
integers in a distinguishable tagged object, preserving the distinction
between an integer and a user string. The decoder restores Python integers
before verification. It does not change Python's global decimal conversion
limit. In particular, masks and multiplicities with more than 16000 bits
are not truncated or passed through floating point.

Proof size is real work. A small number of integer fields can each contain
many bits, and an AHT trace can have many local records. All claims about
certificate cost include the underlying trace lengths, not just the number
of table entries. The algebraic verifier is not a denial-of-service
sandbox; an application must cap serialized input size and provide its
own cooperative deadline when receiving hostile certificates.

\subsection{Two-pass proof construction}
Discovery in the supplied adapter is untraced. A second pass records only
the baseline and physical unions required by \eqref{eq:certificate-queries}.
This avoids retaining every exploratory trace, at the price of recomputing
some counts. The global query and cycle allowances cover both passes. If
proof production fails after successful discovery, the overall certified
call is still \texttt{INCONCLUSIVE}; it does not expose a purported complete
certificate assembled from a prefix.

A coarse upper bound is $Q+V+2$ backend invocations, including the two
baseline counts. Empty unions, repeated masks, and canonical-union reuse
reduce it. A future single-pass producer could retain selected traces or
cache richer verified states, but its memory and replay costs would need
separate analysis. We do not use an unmeasured proof-caching strategy to
improve the reported bound.

\section{Why interval-encoded incidence has polynomial support}
\label{sec:support}

\subsection{The classical weighted route}
For arbitrary abstract nonnegative histograms, $s$ can equal $2^r$.
Interval encoding imposes additional structure. Assign each point the vector
\begin{equation}
 w(x)=(\ind_{A_1}(x),\ldots,\ind_{A_r}(x))\in\NN^r.
 \label{eq:indicator-weight}
\end{equation}
This vector is constant between successive marked endpoints, hence on at
most $2M+1$ integer intervals. The weight of an orbit is the sum of its point
weights, and its nonzero coordinate set is exactly \eqref{eq:signature}.
Thus a compressed list of orbit weights also gives the desired sparse
histogram after equal supports are collected.

We use three classical facts from AHT's analysis \cite{aht}: the classical
scheduler has at most $5k^2(2+\log_2N)$ cycles; weighted transfer preserves
orbit sums and increases the number of live constant-weight intervals by
at most four (Lemma 15 and its proof); and static intervals are archived
as constant-weight interval records rather than individual points.
These are imported ingredients, not new orbit-scheduler proofs.

It already follows from the weighted theorem that the sparse output has
polynomial size. For clarity about the parameter dependence, we give an
explicit, deliberately conservative accounting bound. Tracking both live
and archived runs avoids a possible mistake: a bound on live runs alone
does not automatically bound the cumulative output.

\begin{theorem}[Explicit incidence-support bound]
\label{thm:support}
For an interval system with parameters $(N,k,M,r,B)$ as in
Section~\ref{sec:model}, let
\begin{align}
 L_0&=1+5k^2(2+B),\label{eq:lzero}\\
 S_0&=2M+1+(4k+6)L_0.\label{eq:szero}
\end{align}
Then
\begin{equation}
 s\leq\min\{C,2^r,S_0\}
   =O\bigl(M+(k+1)^3(B+1)\bigr).
 \label{eq:supportbound}
\end{equation}
For $N=0$ the conclusion is immediate. The formula is an upper bound on
occurring incidence types, not on the number of individual orbits.
\end{theorem}
\begin{proof}
The cases $N=0$ and $r=0$ are immediate, so assume both are positive.
Run the classical weighted algorithm conceptually on
\eqref{eq:indicator-weight}. This run is used to prove a structural bound;
the sparse implementation need not execute weighted transfer.
Let $R$ be the number of current constant-weight intervals plus the number
of already archived constant-weight interval records. Initially
$R\leq2M+1$. Normalize equal adjacent live runs after each stage; the archive may
be normalized independently. Temporary endpoint refinement will be charged
explicitly below.

At a static-contraction stage, the union of all pairing domains and ranges
has at most $2k$ interval records, so its complement has at most $2k+1$
static gaps. Refine the current constant-weight partition at the two
endpoints of each gap. This increases the combined run count by at most
$2(2k+1)=4k+2$. Move the pieces contained in gaps from the live list to the
archive and close the corresponding coordinate gaps. Moving an already
refined piece changes its list, not the total number of pieces. Concatenation
or merging decreases, rather than increases, that total.

Identity deletion, reflection trimming, periodic merger, and transmission
leave the weight function unchanged. There is at most one weighted-transfer
operation in a cycle, adding at most four live runs by the imported transfer
bound. The following suffix truncation adds none: it deletes a suffix after
its weight has been transferred. Therefore a complete cycle increases $R$
by at most $(4k+2)+4=4k+6$. The number of pairings never increases, so the
initial $k$ is a valid uniform bound throughout.

The classical cycle estimate, with an extra terminal cycle and with
$B\geq\log_2N$, gives at most $L_0$ cycles. The extra one also covers the
unpaired case $k=0$. At termination every original orbit is represented by
one archived point, but consecutive points with equal vector weight share
an interval record. There are at most $S_0$ such records. Each record
contributes one incidence signature, regardless of its length. Coalescing
records with the same signature cannot increase their number. This proves
$s\leq S_0$. The bounds $s\leq C$ and $s\leq2^r$ are immediate.
\end{proof}

\begin{remark}[What the explicit bound does not say]
The constant in \eqref{eq:lzero} belongs to the classical AHT analysis.
It is not a newly audited cycle constant for the maintained sharp
Fine--Wilf scheduler. The structural conclusion concerns the original
interval relation and is independent of which exact count backend is used
by sparse recovery. The bound is intentionally loose, and no numerical
optimality is claimed.
\end{remark}

\subsection{Consequences and attribution}
\begin{corollary}[No logarithmic port restriction for the local query]
\label{cor:localpoly}
With a polynomial-time exact unweighted interval-orbit backend, the full
sparse incidence profile is computable in time polynomial in
$k+M+r+B$, without assuming $r=O(\log N)$ or
$r=O((\log N)^2)$. With polynomial-size independently replayable count
proofs, the sparse profile also has a polynomial-size independently
checkable certificate.
\end{corollary}
\begin{proof}
Theorem~\ref{thm:support} bounds $s$ polynomially in the encoded parameters.
Theorems~\ref{thm:recovery} and \ref{thm:cert} use at most $O(sr)$ discovery
and verification values. Each coned query has at most $k+2M$ pairings and
the same universe size. Its input, computation, and optional proof therefore
have polynomial size. The remaining operations manipulate polynomially
many $r$-bit masks and $O(B)$-bit counts.
\end{proof}

The polynomial computability of weighted component information predates
this work. The new implementation route is instead a reduction to the
already certified \emph{unweighted} kernel, coupled with a sparse proof
format. It is a strict improvement over the inspected dense interface's
representation bound, not a claim to improve the asymptotic complexity
class of the classical weighted problem.

The full-support static example in Section~\ref{sec:experiments} does not
contradict \eqref{eq:supportbound}. Encoding all $2^r$ signatures using
alternating interval marks requires many marked interval records. What
would be exponentially large as a function of $r$ alone need not be
exponential as a function of the actual interval input length.

\section{Bit costs, storage, and resource semantics}
\label{sec:complexity}
Let $P(k,B)$ be a monotone upper bound for a backend count call. For
certified operation, enlarge $P$ to cover local proof construction, the
proof's serialized bit length, and independent replay. Classical AHT
supplies a polynomial count bound; a chosen implementation must separately
meet the corresponding proof contract.

Define $D_s=\sum_{T\in\supp h}|T|$ and let $Q,V$ be the bounds in
\eqref{eq:qsum} and \eqref{eq:vbound}. A conservative bit-time bound for
our direct subset-scanning implementation is
\begin{equation}
\begin{split}
 &(Q+V+2)P(k+2M,B)\\
 &\quad+O\left((Q+V+s+1)
 \left[r^2+s(B+r+1)
 +(M+1)\log(M+2)(B+r+1)\right]\right),
\end{split}
\label{eq:bitcost}
\end{equation}
up to polynomial input-validation and proof-serialization overhead.
The displayed dictionary-cache accounting uses expected hash-table costs.
Replacing it by a deterministic ordered dictionary adds comparison and
logarithmic factors and preserves the polynomial conclusion. A simple
linear scan of the polynomial-size cache also gives a deterministic
polynomial bound without a hashing assumption.

Here is the accounting behind \eqref{eq:bitcost}. An ordinary implementation
scans port masks using $r$-bit arithmetic, safely bounded by $O(r^2)$ per
query. A residual value checks at most $s$ subset relations and subtracts
$O(B)$-bit multiplicities, costing $O(s(r+B+1))$. Forming a selected physical
union sorts at most $M$ interval records and compares $B$-bit endpoints.
The verifier's one-bit deletions and sparse table checks fit the same
coarse bound. These operations are not assumed to be unit-cost arithmetic
on arbitrarily large integers.

Ignoring backend traces, the sparse profile occupies $O(s(r+B+1))$ bits.
Cached physical unions occupy at most $O((Q+V)M(B+1))$ bits in a straightforward
representation. A two-pass certificate holds at most $V+1$ distinct count
proofs before physical-union reuse, together with its sparse index table.
There is no hidden output array of length $2^r$.

The result is polynomial after substituting \eqref{eq:supportbound}, but
it is not necessarily fast for all moderate-size inputs. The simple
residual implementation can cost quadratic time in $s$ across all probes.
The deliberately loose structural bound can also be much larger than
observed support. Output-sensitive accounting remains useful precisely
because it does not replace actual $s$ by that worst-case upper bound.

\subsection{Three separate budgets}
The implementation distinguishes support entries, backend invocations,
and AHT cycles. A cycle allowance counts structural iterations, not all
bit operations or elapsed time. A \texttt{check()} callback is forwarded
to the maintained backend and polled during the sparse search. Input
normalization and individual Python operations are not a hard real-time
preemptive boundary, so a separate process-level policy is appropriate
for adversarial inputs.

For a fixed encoded instance with all caps disabled, the preceding theorems
apply. With a cap, no incomplete output is promoted to a complete profile.
The native adapter shares its backend query and cycle allowances between
ordinary discovery and optional proof production. Exhaustion publishes only
status, reason, and work statistics. No local status is converted to
\texttt{UNKNOT} or \texttt{KNOTTED}.

\section{Incidence-resolved binary orientation data}
\label{sec:orientation}
An interval's order reversal is not the same datum as a surface gluing's
orientation character. Let each supplied pairing $g$ carry a separate
parity $\epsilon_g\in\{0,1\}$. On $X_N\times\{0,1\}$ form the relation
\begin{equation}
 (x,e)\sim(g(x),e+\epsilon_g\bmod2).
 \label{eq:cover}
\end{equation}
This is encoded by $2k$ interval pairings on two consecutive blocks of
length $N$. Lift every port to both sheets. There are at most $2M$ marked
interval records and still $r$ named ports.

A base component is \emph{consistent} if its vertices admit binary labels
whose differences equal the supplied edge parities. Equivalently, every
closed walk has parity zero. This is a statement about the supplied signed
relation. It becomes geometric orientability only when the parities have
been derived and verified as the relevant orientation character.

\begin{theorem}[Support-preserving orientation refinement]
\label{thm:orientation}
Let $h$ and $\widetilde h$ be the base and lifted incidence histograms.
Their supports are identical. If $o(T)$ and $u(T)$ are the numbers of
consistent and inconsistent base components of signature $T$, then
\begin{align}
 \widetilde h(T)&=2o(T)+u(T),&h(T)&=o(T)+u(T),\label{eq:covercounts}\\
 o(T)&=\widetilde h(T)-h(T),&u(T)&=2h(T)-\widetilde h(T).\label{eq:orientresult}
\end{align}
Once $h$ is known, all coefficients of $\widetilde h$ can be recovered using
at most $s$ additional lifted zeta values, plus one lifted baseline count.
\end{theorem}
\begin{proof}
A path in a base orbit lifts uniquely once its initial sheet is chosen.
Every component upstairs therefore projects onto the entire base orbit.
If all closed walks have parity zero, the two choices of initial sheet
remain disconnected and give two lifted components. If some closed walk
has parity one, the two sheets are connected, giving one lifted component.

Every lifted component projects onto every point of its base orbit and the
ports include both sheets. Thus every lift has exactly the same port
signature as the base orbit. No new signature is introduced and none is
lost. Counting one or two lifts gives \eqref{eq:covercounts} and solving
the two equations gives \eqref{eq:orientresult}.

Order the known support by nondecreasing cardinality, breaking ties
arbitrarily. A proper subset always precedes its superset. For each known
$T$ compute
\begin{equation}
 \widetilde h(T)=\widetilde G(T)
 -\sum_{\substack{S\in\supp h\\S\subsetneq T}}\widetilde h(S).
 \label{eq:triangular-cover}
\end{equation}
Support equality proves there are no unlisted terms. One lifted zeta value
per type suffices. The baseline checks the total; the necessary inequalities
$h(T)\leq\widetilde h(T)\leq2h(T)$ are useful defensive checks, not a
replacement for the exact cover construction.
\end{proof}

The generic cover construction and known-support calculation are implemented
and tested. A native source-bound certificate wrapper for the combined
signed result is not included: it must authenticate the signed input,
both-sheet port lift, and coned cover counts in addition to the base sparse
certificate. The theorem shows that this wrapper needs no second support
search. It does not license treating unauthenticated parity bits as a
surface-orientability certificate.

\begin{example}[A reflection with a fixed point]
On five points, pair $x$ with $4-x$ and assign parity one. The two two-point
orbits are consistent; the midpoint has a parity-one loop and is inconsistent.
With one port containing every point, $h(\{1\})=3$ and
$\widetilde h(\{1\})=5$. Formula~\eqref{eq:orientresult} gives two consistent
and one inconsistent component. This case tests why a geometric reflection
flag must not be substituted for a consistency conclusion.
\end{example}

\section{Exact continuation operations and a sharp limitation}
\label{sec:continuations}

\subsection{Coning without another orbit calculation}
Once the complete profile is known, some later quotient operations can be
performed directly on it. For $J\subseteq[r]$, cone all points of
$\bigcup_{i\in J}A_i$ in the current relation.

\begin{theorem}[Coning update]
\label{thm:coneupdate}
Let $\mathcal T_J=\{T\in\supp h:T\cap J\neq\varnothing\}$.
If $\mathcal T_J$ is empty, the profile is unchanged. Otherwise remove
all entries of $\mathcal T_J$ and insert the single entry
\begin{equation}
 \left(\bigcup_{T\in\mathcal T_J}T,\ 1\right).
 \label{eq:coneupdate}
\end{equation}
All other entries keep their multiplicities. The update is exact and does
not increase the support size.
\end{theorem}
\begin{proof}
Every current component meeting the selected union is joined to every
other such component. Its new port signature is the union of their old
signatures. This describes exactly one new component even if a selected
signature had enormous multiplicity. Untouched components remain separate
and retain their signatures. The new signature intersects $J$, so it cannot
coincide with an untouched signature. Replacing one or more selected types
by one therefore does not increase support size.
\end{proof}

A sequence of $q$ such operations costs $O(qs)$ mask and count operations
in a direct scan, or $O(qs(r+B+1))$ conservative bit work with fixed ports.
The same theorem applies inductively after each quotient. No fresh AHT call
is required for this restricted continuation.

\subsection{Union re-marking and disjoint unions}
Suppose each new port is a union of old ports,
$B_j=\bigcup_{i\in J_j}A_i$. Map a signature by
\begin{equation}
 \phi(T)=\{j:T\cap J_j\neq\varnothing\}.
 \label{eq:relabel}
\end{equation}
The new profile is obtained by adding the multiplicities of types with
equal $\phi(T)$. This follows directly from the existential definition of
port incidence. Deleting or renaming ports is a special case. If two
independent systems are combined disjointly and each named port is the
disjoint union of its two old versions, their profiles add coefficientwise.

\begin{corollary}[Restricted continuation congruence]
Two systems with the same incidence profile remain indistinguishable by
any finite composition of coning selected port unions, union re-marking,
and disjoint union with fixed auxiliary systems, when the observation is
the resulting incidence profile.
\end{corollary}
\begin{proof}
Each operation has a well-defined update depending only on the current
profile. Induct on the number of operations.
\end{proof}

This is a sufficient-state theorem for an explicit operation language,
not a claim that incidence is a complete boundary state for a topological
hierarchy. The distinction is essential for subsequent integration.

\subsection{Why arbitrary attachments are not covered}
\begin{proposition}[Equal histograms can separate under an attachment]
\label{prop:attachment}
There exist two interval-generated equivalence relations with identical
port-incidence profiles for which adding the same singleton pairing gives
different component counts.
\end{proposition}
\begin{proof}
Use four points and a single port containing all four. The first relation
has classes $\{0,1\},\{2,3\}$; the second has classes
$\{0,2\},\{1,3\}$. Both have histogram $h(\{1\})=2$ and all other entries
zero. Both relations are generated by singleton interval pairings. Adding
$0\sim1$ changes nothing in the first system, but joins the two classes in
the second. Their resulting component counts are two and one respectively.
\end{proof}

An incidence histogram also cannot answer pointwise-intersection questions
by substituting logical conjunction of ports. A component can meet $A$ and
$B$ at different points even when $A\cap B=\varnothing$. More generally,
it contains no point multiplicities, cyclic attachment orders, slopes,
component Euler characteristics, embeddings, or individual attachment maps.
Some of these can be obtained from richer weighted or geometric data;
none can be silently reconstructed from a Boolean signature alone.

\section{Implementation, trust boundaries, and integration}
\label{sec:implementation}

\subsection{Package modules}
The artifact uses only the Python standard library. The generic
\texttt{sparse\_ports/core.py} implements both deletion strategies, sparse
certificate payloads, independent algebraic checking, coning, and union
re-marking. \texttt{intervals.py} implements normalized physical unions,
the coning oracle, and a bounded dense reference for controlled comparisons.
\texttt{orientation.py} supplies the double cover and known-support
refinement. \texttt{codec.py} handles tagged hexadecimal transport.

The opt-in \texttt{proveit.py} adapter imports the maintained orbit modules
lazily, so the standalone research tests need no installation of ProveIt.
It preserves the maintained convention that incomplete calculations expose
no certified count. Its sparse histogram has an explicit encoding tag;
a caller expecting the old dense array cannot accidentally treat a sparse
pair list as the same API.

For example, with the maintained \texttt{fast} directory and this package
on \texttt{PYTHONPATH}, a caller can use:
\noindent\begin{minipage}{\linewidth}
\begin{lstlisting}[language=Python]
from fastunknot.interval_orbits import IntervalPairing
from sparse_ports.proveit import (
    analyze_port_incidence_sparse,
    verify_sparse_port_incidence_certificate,
)
pairs = [IntervalPairing(0, 3, 4, 7)]
ports = [[(0, 2)], [(1, 3)], [(7, 10)], []]
answer = analyze_port_incidence_sparse(
    10, pairs, ports, record_certificate=True,
)
if answer["status"] == "COMPLETE":
    assert verify_sparse_port_incidence_certificate(
        10, pairs, ports, answer["certificate"],
    )
\end{lstlisting}
\end{minipage}
This illustrates the intended integration call, not a claim that this
exact native call was executed in the delivery environment.

\subsection{The audit extraction is not the maintained implementation}
The executed component experiments use \texttt{tests/aht\_reference.py},
an explicitly labeled audit extraction of the inspected maintained
unweighted scheduler. It keeps the interval reductions and sharp merger,
but omits native proof recording and unrelated wrappers. It is not a
byte-identical repository snapshot. Source attribution and the inspected
blob identifier are retained in its header.

Small systems are checked against an independent point-expanded disjoint-set
calculation. Huge static and periodic systems are checked against a
non-expanding endpoint/residue calculation. The dense and sparse benchmark
arms use exactly the same audit backend and physical-union cache, isolating
the reconstruction strategy. The dense arm is a functional reference with
the same dense-inversion architecture, not a byte-for-byte run of the
maintained dense module.

Consequently these experiments support a local algorithmic comparison.
They do not measure the maintained proof producer, certify its current
installation, establish native normal-surface provenance, or show a
whole-knot speedup. A claim of such a result would require a different
measurement and validation boundary.

\subsection{Proposed non-disruptive integration}
Keep this work initially under a research subdirectory and expose the
sparse interface alongside, not in place of, the dense function. Run
\texttt{integration/native\_smoke.py} in a real checkout. That script compares
the current dense and sparse profiles on randomized systems, records and
replays actual maintained local proofs, round-trips the certificate codec,
and disables the producer while replaying. It is supplied but unexecuted
here.

After those checks, add native normal-arc fixtures and paired query
measurements. Only then consider promoting an opt-in sparse backend in the
maintained module. Dense output should remain an explicit conversion with
an explicit port limit, because producing every zero entry is intrinsically
an output-size operation. Nothing in this package changes ordinary
\texttt{recognize} dispatch or removes its complete fallback.

\section{Validation and controlled component experiments}
\label{sec:experiments}

\subsection{Executed checks}
The included test log records 19 passing unittest methods. A method often
checks many instances: the total is not merely 19 examples. The executed
coverage includes all 6654 histograms with $r\in\{0,1,2,3\}$ and coefficients
in $\{0,1,2\}$, comparing both discovery strategies and certificate replay;
1000 random sparse abstract profiles; and 80 tests of the balanced-block
bound on singleton supports. Large weights, empty ports, zero universes,
malformed integers, callbacks, caps, and mutated certificates are included.

Five hundred random interval systems compare the audit AHT count with
literal components and compare sparse, dense, and literal histograms.
Three hundred further systems undergo five successive conings each,
checking the profile update after every operation against literal quotient
connectivity. Huge periodic cases use endpoint lengths up to roughly 4000
bits, and huge static and transport cases reach 16001-bit integers without
point expansion. Five hundred random signed systems compare the refined
profile with an independent breadth-first parity-consistency calculation.

The maintained-adapter contract tests inject a backend with independent
literal count witnesses. They test orchestration, exact proof bindings,
shared budgets, codec round trips, and replay with the producer disabled.
They are explicitly \emph{not} execution of the maintained AHT trace
producer/verifier. No statement about the repository's full test count or
Regina integration is inferred from this local test run.

These checks are evidence for implementation agreement, not substitutes for
the proofs. In particular, finite tests cannot establish the structural
support bound for every interval system. The latter depends on the
classical transfer and cycle facts plus the accounting in
Theorem~\ref{thm:support}.

\subsection{Measurement protocol}
The retained driver uses CPython 3.13.5 on the Linux platform recorded in
\texttt{results/benchmarks.json}. The fixed shuffle seed is 261008505.
Nine paired workloads have four arms: identical dense controls A and B,
balanced splitting, and individual deletion. Each workload receives one
warmup round and seven measured rounds, with arm order shuffled per round.
Three larger capacity workloads run only the sparse arm. This gives 273
measured calls and 39 warmup calls. Correctness comparison is outside the
timed region. Raw nanoseconds, query counts, all arms, and source hashes are
retained, not just the favorable ratios.

The reported ratio is the median of seven within-round dense-A/sparse
ratios, not the ratio of independently selected fastest runs. A value above
one favors sparse recovery. Timing includes oracle creation and its baseline
count, union preparation, all backend calls, and reconstruction. It excludes
native local proof production and replay, which this benchmark does not
execute. Timings are local measurements, not extrapolations to other
hardware or an unknot corpus.

\begin{table}[t]
\centering\small
\begin{tabular}{lrrrrr}
\toprule
Workload & $s$ & Dense A (ms) & Dense B (ms) & Split (ms) & Paired ratio\\
\midrule
coincident-8 & 2 & 0.586 & 0.591 & 0.052 & 11.00\\
disjoint-8 & 9 & 28.011 & 27.187 & 3.682 & 7.40\\
coincident-10 & 2 & 2.824 & 2.838 & 0.070 & 40.17\\
disjoint-10 & 11 & 176.143 & 173.715 & 7.946 & 22.08\\
coincident-12 & 2 & 15.264 & 14.052 & 0.120 & 122.96\\
disjoint-12 & 13 & 1003.604 & 995.129 & 14.636 & 68.04\\
periodic-6 & 7 & 1.315 & 1.263 & 0.556 & 2.45\\
periodic-8 & 9 & 6.139 & 5.862 & 0.861 & 6.66\\
full-support-6 & 64 & 14.773 & 14.859 & 20.208 & 0.73\\
\bottomrule
\end{tabular}

\caption{Controlled component timings. Dense A and B are identical controls.
The last column is the paired dense-A/split ratio; values below one are
regressions. No native proof-generation or whole-knot time is included.}
\label{tab:timings}
\end{table}

\begin{table}[t]
\centering\small
\begin{tabular}{lrrrr}
\toprule
Workload & Dense calls & Split calls & Split requests & Linear requests\\
\midrule
coincident-8 & 2 & 2 & 16 & 16\\
disjoint-8 & 256 & 21 & 57 & 72\\
coincident-10 & 2 & 2 & 20 & 20\\
disjoint-10 & 1024 & 26 & 79 & 110\\
coincident-12 & 2 & 2 & 24 & 24\\
disjoint-12 & 4096 & 33 & 101 & 156\\
periodic-6 & 37 & 15 & 49 & 42\\
periodic-8 & 114 & 21 & 71 & 72\\
full-support-6 & 64 & 64 & 606 & 384\\
\bottomrule
\end{tabular}

\caption{Backend calls and logical discovery requests on the same workloads.
The dense reference evaluates all subset masks, but equal physical unions
share backend work. Split requests can exceed distinct backend calls.}
\label{tab:queries}
\end{table}

\subsection{What the families test}
For the coincident family, $N=2^{500}$, there are no base pairings, and all
$r$ ports equal $[5,N-7)$. The only signatures are the empty one, with
multiplicity 12, and $[r]$, with multiplicity $N-12$. Both dense and sparse
methods use only two backend calls because of physical-union reuse. Their
different time costs therefore isolate the dense mask/inversion overhead
rather than faster orbit counting.

For the disjoint family, $r$ separated two-point intervals mark otherwise
static points. Its $r+1$ occurring types are the empty signature and the
$r$ singleton signatures. Distinct port subsets have distinct physical
unions, so the dense method cannot avoid exponentially many coned calls by
union caching. Balanced extraction removes most absent combinations and
also uses fewer logical probes than individual deletion on this family.

For the periodic family, a single pairing translates by
$p=2^{250}+37$ on a universe of size $N=p(2^{250}+1)$. Components are residue
classes modulo $p$. Overlapping interval ports are checked independently
by sweeping their residue endpoints. This family exercises actual compressed
orbit reduction rather than a purely static universe, but it is still a
controlled interval system, not an extracted knot hierarchy.

The full-support control uses $r=6$, $N=64$, no base pairings, and interval
ports that encode the six binary digits of each point. All 64 signatures
occur. Its interval encoding has 63 marked interval records; sparsity offers
no output reduction. Balanced recovery is slower than dense inversion,
approximately 20.2 ms versus 14.8 ms in the retained medians. The paired
ratio is about 0.73. This negative control rules out a claim that the sparse
method is uniformly faster.

On the sparse workloads, paired ratios range from approximately 2.45 to
123 in the retained data. These figures quantify this particular
component-level ablation. They do not compare a fully optimized weighted
AHT implementation against sparse coning, and do not show a speedup in
ordinary recognition. Identical-control variability is visible in the raw
samples; precision beyond the displayed few significant digits would not
be justified.

\subsection{Large-port capacity without a dense comparison}
The capacity family uses $N=2^{4000}$ and the same two-type coincident-port
structure. Table~\ref{tab:capacity} reports completed sparse calculations.
No dense timing is attempted at 32, 128, or 1024 ports, and no speed ratio
is fabricated for those omitted runs.
\begin{table}[h]
\centering\small
\begin{tabular}{rrrrr}
\toprule
Ports $r$ & Endpoint bits & Support $s$ & Orbit calls & Median (ms)\\
\midrule
32 & 4001 & 2 & 2 & 0.298\\
128 & 4001 & 2 & 2 & 2.635\\
1024 & 4001 & 2 & 2 & 151.524\\
\bottomrule
\end{tabular}

\caption{Completed large-port sparse queries. These are capacity observations,
not ratios against an executed dense arm. Endpoint size is binary length.}
\label{tab:capacity}
\end{table}

Only two backend count calls are needed at each size, but elapsed time still
increases with $r$: constructing and comparing large masks and repeatedly
forming unions is not free. Generic algebraic certificate-query accounting
is also retained in the JSON, but is separate from these timed calls and
from genuine local AHT trace costs. Quotienting provably identical ports
before discovery could reduce this overhead, subject to a source-bound
expansion map for the output and certificate.

\section{Consequences for the research program}
\label{sec:global}

\subsection{A local exponential factor is avoidable}
The inspected dense representation spends $2^r$ time and output slots even
when interval encoding permits only polynomially many occurring signatures.
The sparse interface removes that representation penalty using the existing
unweighted count boundary. Theorem~\ref{thm:support} is stronger than an
assumption that a particular benchmark happens to have small support:
it makes the local polynomial statement uniform over interval-encoded
instances.

This conclusion should not be advertised as an improvement from exponential
to polynomial \emph{unknot recognition}. Its unit of input is a supplied
pairing system and marked interval list. A complete recognition procedure
must still construct those systems, justify what they mean geometrically,
choose successful decompositions, and bound all intermediate description
sizes and the number of calls.

\begin{proposition}[Conditional composition, with explicit obligations]
Suppose a sound and complete recognition procedure on an $n$-crossing input
makes at most $n^{O((\log n)^a)}$ local incidence calls. Suppose every call
has $k+M+r+B\leq n^{O((\log n)^b)}$, its geometric provenance is verified
within the same kind of bound, and all other work has quasi-polynomial bit
cost. Replacing dense incidence with the sparse exact interface contributes
only quasi-polynomial time to that procedure.
\end{proposition}
\begin{proof}
By Corollary~\ref{cor:localpoly}, each call costs a fixed polynomial in its
encoded parameter sum, hence at most $n^{O((\log n)^b)}$. Multiplying by
the call bound remains quasi-polynomial. Add the separately bounded
provenance and non-incidence work. Correctness requires the stated soundness
and completeness hypotheses; it is not supplied by the time calculation.
\end{proof}

The point is not that this elementary composition argument proves a new
global complexity theorem. It identifies exactly which assumptions remain
outside the local interface. The logarithmic-port hypothesis is no longer
one of them. Bounds on $M$ and $B$ after repeated cutting, and a sufficient
state for actual attachment maps, remain indispensable.

\subsection{A precise correction to an interpretation of the old bound}
The earlier $2^r$ estimate is correct for the maintained dense algorithm.
It would be incorrect to interpret it as an intrinsic lower bound for
computing all \emph{occurring} interval-incidence types. The weighted AHT
route already prevents that interpretation; the present sparse reduction
makes the distinction executable with the current unweighted proof interface.
No change to the correctness of the old dense routine is alleged.

Likewise, equal signatures may justify batching certain aggregate queries,
but they cannot generally justify merging components in a geometric state.
Proposition~\ref{prop:attachment} shows the missing information explicitly.
Any future integration that discards attachment maps merely because two
components have the same signature would need a stronger theorem than the
one proved here.

\section{Further research questions and concrete next experiments}
\label{sec:future}
The following questions are ordered roughly from immediate integration work
to the broader hierarchy problem. Each asks for an additional result or
measurement rather than treating a missing hypothesis as established.

\paragraph{1. Genuine native proof replay and differential testing.}
Run the supplied native integration script at a pinned maintained commit,
then add adversarial mutations of every sparse proof binding and repeat
against both classical and sharp merger proof versions. Does the same
independent sparse verifier accept both without importing either scheduler?
Retain the actual source hashes and the full native test log.

\paragraph{2. Native normal-arc and disc fixtures.}
Extract ports from actual validated normal-arc systems, rather than choosing
coordinate intervals artificially. Measure the observed $(k,M,r,B,s)$ and
the distribution of signature sizes. Which geometric constructions produce
small $s$ and small $D_s$, and which fragment the port representation?
A geometric port extractor should have its own independently replayable
provenance contract.

\paragraph{3. A signed certificate adapter.}
Implement the source-bound double-cover wrapper suggested by
Theorem~\ref{thm:orientation}. Prove and test that the checker binds every
parity and every lifted port to the source, then compare known-support
refinement against an independent second sparse search. The generic
savings of one support search should be separated from actual trace costs.

\paragraph{4. One-pass weighted AHT versus sparse coning.}
Build or adapt a verified weighted transfer kernel and compare it fairly
against this black-box reduction. The weighted algorithm can reuse more
internal structure; sparse coning reuses an existing smaller trusted kernel.
Which method minimizes bit operations, proof bytes, and replay time on
actual normal-surface families? A hybrid should switch based on measured
work, not merely on a port-count threshold.

\paragraph{5. Quotient identical or implication-related ports.}
Identical physical ports are safe to quotient with an explicit output
expansion map. More generally, geometric containment $A_i\subseteq A_j$
forces every signature containing $i$ to contain $j$. Can a canonical
implication-poset representation reduce both search and certificate probes
while preserving independent source checks? Numerical equality of coverage
counts alone is not a sufficient justification.

\paragraph{6. Faster residual subset sums.}
The prototype scans all previously found types for each residual query.
Develop a deterministic incremental data structure for
$\sum_{T\subseteq U}h(T)$ restricted to the discovered support, with bit-cost
bounds that improve this quadratic factor on geometrically occurring
families. A generic dense zeta table would reintroduce the very cost being
removed.

\paragraph{7. Sharpen the structural support estimate.}
The conservative bound $O(M+(k+1)^3(B+1))$ charges every possible static
gap in every cycle. Can live/archive amortization remove a factor of $k$
or weaken the dependence on $B$? Determine matching lower-bound families
for the number of distinct signatures, not merely for the number of
components or distinct orbit weight vectors.

\paragraph{8. Optimal multi-support query complexity.}
The singleton lower bound does not settle the optimal reconstruction cost
for arbitrary positive profiles. Characterize the best adaptive additive
query bound in terms of $r$, $s$, signature sizes, and inclusion structure.
Can a shared extraction tree avoid repeatedly discovering common mandatory
ports? Any improvement must account for multiplicities rather than
assuming unweighted support.

\paragraph{9. A richer exact continuation state.}
Find the smallest useful extension of the histogram that is sufficient for
specific attachment maps arising in compressed cutting. One candidate is
a sparse incidence profile together with a compressed boundary permutation
or interval map. Prove a congruence theorem for the actual operation language,
and exhibit separating examples when proposed state fields are omitted.

\paragraph{10. Port fragmentation under cutting.}
Even a polynomial local algorithm can be overwhelmed by an exponentially
fragmented input. Establish bounds on the number of port intervals after
one certified cutting operation and under a hierarchy. Can periodic or
straight-line encodings replace long interval-union lists without destroying
the support bound or the coning reduction?

\paragraph{11. Global hierarchy costs and completeness.}
Relate the encoded sizes of actual intermediate objects to the hierarchy
length and pattern-complexity parameters. The target is a constructive
bound on total bit work, including successful search, not just fast checking
of a supplied hierarchy or its aggregate orbit data. This is the missing
connection to a general quasi-polynomial recognizer.

\paragraph{12. Whole-pipeline evaluation and proof formalization.}
After native integration, measure complete recognition calls with paired
controls, including adverse cases and certificate replay. Separately,
formalize the nonnegative-residual and minimal-support certificate theorems;
they form a compact first target because they do not require formalizing
three-manifold topology. Only promote default behavior when correctness
and full-call performance have both been established.

\appendix
\section{A worked sparse certificate}
\label{app:example}
Consider three ports and the profile
\[
 h(\varnothing)=2,\quad h(\{1\})=3,\quad
 h(\{2,3\})=5,\quad h(\{1,2,3\})=7,
\]
with all other coefficients zero. The baseline is $C=17$. The order
$\varnothing,\{1\},\{2,3\},\{1,2,3\}$ is a valid minimal-extraction order.
For the empty type, the single check is $G(\varnothing)=2$. For $\{1\}$,
$G(\{1\})=5$ and subtracting the earlier weight two gives three; deleting
port one leaves residual zero.

For $\{2,3\}$, $G(\{2,3\})=7$, while
$G(\{2\})=G(\{3\})=2$. Subtracting the previously recovered empty type
leaves five at the pair and zero at both deletions. For the full type,
$G(\{1,2,3\})=17$; the three earlier weights sum to ten. Its deletion values
are $G(\{2,3\})=7$ and
$G(\{1,2\})=G(\{1,3\})=5$, exactly exhausted by earlier contained types.
The residual at the full type is seven and each deletion residual is zero.
Finally $2+3+5+7=17$ proves no omitted positive coefficient.

This example distinguishes the independent certificate from a log of the
search. The verifier does not need to know which balanced blocks were
visited, which candidate unions were cached, or which failed probes the
producer made before finding these types. It checks only the local
minimality conditions and total mass.

\section{Reproduction and artifact inventory}
\label{app:reproduce}
From the package root, the standalone tests and retained-data tables are
reproducible with:
\begin{lstlisting}[language=bash]
python -m unittest discover -s tests -v
python benchmarks/make_tables.py
cd article
sh build.sh
\end{lstlisting}
To obtain new timing samples without overwriting the retained experiment:
\begin{lstlisting}[language=bash]
python benchmarks/run_benchmarks.py \
  --output results/benchmarks_rerun.json
\end{lstlisting}
The driver and article-table generator deliberately use different filenames
for an unreviewed rerun and the retained evidence. To execute genuine native
integration in an available checkout:
\begin{lstlisting}[language=bash]
python integration/native_smoke.py /path/to/ProveIt --cases 250
\end{lstlisting}
This last command is a required next validation step, not a completed
experiment in the supplied test log.

The archive contains the TeX source, compiled PDF, generated table fragments,
all research modules and tests, the audit-only kernel with attribution,
raw benchmark data, test log, source provenance, a claim-status ledger,
build/reproduction scripts, and integration guidance. It contains no
redistributed third-party research PDFs or font files. Its code is offered
under MIT-0 for compatibility with the inspected project.

\clearpage
\begin{thebibliography}{9}
\raggedright
\bibitem{aht}
I. Agol, J. Hass, and W. Thurston,
\emph{The computational complexity of knot genus and spanning area},
Transactions of the American Mathematical Society \textbf{358} (2006),
3821--3850. See Theorem 12, Section 6, Lemma 15, and Theorem 16.
\href{https://arxiv.org/abs/math/0205057}{arXiv:math/0205057}.

\bibitem{lackenby2026}
M. Lackenby,
\emph{Incompressible surfaces, hierarchies and unknot recognition},
2026 preprint, version 1. See Section 9.
\href{https://arxiv.org/abs/2607.23350}{arXiv:2607.23350}.

\bibitem{bshouty-mazzawi}
N. H. Bshouty and H. Mazzawi,
\emph{Optimal query complexity for reconstructing hypergraphs},
2010 preprint.
\href{https://arxiv.org/abs/1001.0405}{arXiv:1001.0405}.

\bibitem{price-scarlett}
E. Price and J. Scarlett,
\emph{A fast binary splitting approach to non-adaptive group testing},
APPROX/RANDOM 2020, Leibniz International Proceedings in Informatics,
Article 13.
\href{https://doi.org/10.4230/LIPIcs.APPROX/RANDOM.2020.13}{doi:10.4230/LIPIcs.APPROX/RANDOM.2020.13}.
\href{https://arxiv.org/abs/2006.10268}{arXiv:2006.10268}.

\bibitem{repo-incidence}
V. Reshetnikov, ProveIt repository,
\emph{UnknotRecognition: interval incidence and independently replayable
interval orbits}. Inspected modules
\texttt{fast/fastunknot/interval\_incidence.py},
\texttt{interval\_orbits.py}, and \texttt{interval\_orbit\_verify.py}.
Incidence integration baseline: commit
\texttt{0441aad3967b69621edabfdb85367d465b4fb32f} (8 October 2026).
\href{https://github.com/VladimirReshetnikov/ProveIt/tree/0441aad3967b69621edabfdb85367d465b4fb32f/Topology/UnknotRecognition}{Pinned source directory}.
Individual inspected blob hashes are in \texttt{provenance.json}.

\bibitem{repo-article}
ProveIt, \emph{Independent orbit proofs and reusable marked-union queries},
\texttt{synthesis/orbit\_certificates.tex}, inspected blob
\texttt{a83b3ab44d95b2e8c35752266d9db626ca126225}.
\href{https://github.com/VladimirReshetnikov/ProveIt/blob/main/Topology/UnknotRecognition/synthesis/orbit_certificates.tex}{Maintained article section}.
Accessed 8 October 2026; the moving link is not a claim of immutable content.
\end{thebibliography}
\end{document}
